\section{Bonfires, nodes and receipts}\label{sec:encounters}

\paragraph{Participants.} Two kinds of party take part, on a host chain that the protocol takes as given.
\begin{itemize}
  \item A \emph{Bonfire} is a community with a grammar. Anyone who takes part in its discourse is a \emph{member},
  at no cost; in the ordinary case a member is a person talking to the Bonfire's agent. A \emph{stakeholder} is a
  member who holds stake in the Bonfire, either earned through contribution (Section~\ref{sec:reward}) or locked
  from outside. Stake weighs in deciding what enters the grammar, shares in what the Bonfire earns, and is what is
  lost when the Bonfire is at fault (Section~\ref{sec:bond}). A newcomer to the network joins a Bonfire that
  exists, or starts one and takes up constructions from the others.
  \item A \emph{node} supplies computation. It holds the inventory of each Bonfire it serves (Section~\ref{sec:constructions}), runs the engine on what
  is said, issues receipts, and re-executes the receipts of other nodes.
  To serve a Bonfire a node commits a fixed \emph{pledge} to it (Section~\ref{sec:bond}). A node may be run by
  anyone, including the Bonfire it serves.
\end{itemize}
The Bonfire that registered a construction is its \emph{originator}. Two Bonfires are \emph{independent} when no
party controls the governing stake of both. Members in common do not make two Bonfires one: a person who belongs
to two communities is the ordinary route by which a construction travels from one to the other.

\paragraph{The encounter.} The elementary act of the protocol is an \emph{encounter}: one round of a naming game
\citep{steels1995spatial,baronchelli2006sharp}, in which a speaker produces a form for a meaning and a hearer
obtains a meaning from it. Here at least one party is an agent applying an inventory, and every encounter
takes place in the discourse of some Bonfire. In the ordinary case the other party is a member, talking to an agent
that may direct sub-agents of its own. The agents of two Bonfires may also meet.

\paragraph{What can be replayed.} When an agent reads what was said to it, the engine applies constructions from
the inventory to the text and either arrives at a complete reading or does not. The protocol requires that step to
be deterministic: anyone who holds the text and the inventory can repeat it and obtain the same reading, and can
remove one construction and see whether the reading survives. The protocol pays only on facts of this kind.
Whether the person was satisfied, or whether an outside system accepted a step the agent took, is not replayable.
Such outcomes move a construction's standing in the inventory (Section~\ref{sec:constructions}), and they are no
part of what is paid.

Determinism is a requirement on the implementation and is listed among the open problems
(Section~\ref{sec:open}). The constructions in a receipt are selected by the engine's search with no model in the
loop, and the receipt fixes by hash the engine version, the inventory it was read against, and the order in which
constructions are tried. Where a model settled a choice the grammar left open, the choice is recorded in the
receipt, or the use is not counted.

\paragraph{The receipt.} Each encounter produces a \emph{receipt}: a record naming the constructions that were
necessary to the reading, by their hashes, with a commitment to the text, the Bonfire in whose discourse the
encounter took place, the node, and the epoch. The node signs it. The receipt is what the network reasons about;
the conversation itself stays where it happened.

Nothing about a receipt needs a ledger until it is disputed. Receipts are therefore issued as signed attestations
held off the chain and committed to in bundles, a pattern of which the Ethereum Attestation Service is one existing
realisation \citep{eas2024contracts}. A receipt is signed as structured data, so anyone can verify the signature
without consulting a ledger. Each receipt references the one before it in the same conversation, and a
registration references its parents, so a conversation and a lineage are both chains of attestations. Once per
epoch a node records on the host chain a single hash, the root of a tree over all its receipts for the epoch, which
commits it to them before credit is computed. After a period in which any node may challenge, a contract on the
host chain settles the epoch from the committed receipts.

Three rules follow. Receipts are irrevocable, so that a node cannot retract one after it has been counted. A
receipt under a committed root that its node does not produce on challenge is treated as one that failed to
re-execute. And the text is attested only as the root of a tree over its turns, so that a challenge discloses the
turns of that one encounter to the challenger and nothing else. That disclosure is a cost of the design
(Section~\ref{sec:open}). An attestation establishes who signed what, and no later than when. That what was signed
is true is established by re-execution and by stake.

\paragraph{Gaps.} When an agent cannot resolve a form, the node records the failure as a \emph{gap}: a form, in a
context, for which the inventory had no reading. Gaps are a Bonfire's own record of what its members tried to say
and could not. They are the network's demand, and they are where the taking up of a construction begins
(Section~\ref{sec:adoption}).
